\section{统一构造：精确积分加残差抽样}
\label{sec:kl-residual-framework}

KL 数值、KL 梯度和 baseline 的统计分析可以共用一个积分问题。固定当前 prefix 后，$p$、$q$、$r_v$ 与 $s_v$ 都是词表上的确定函数。选择积分测度 $\nu$ 和被积函数 $f$，便确定了估计目标；选择实际抽样分布 $\rho$ 和可精确积分的近似函数 $a$，便确定了估计方法。以下所有协方差均针对动作抽样，向量的方差仍指协方差迹。

\begin{definition}[可精确积分的近似函数与残差]
\label{def:kl-residual-integral}
令 $\nu\in\{p,q\}$，$f:\mathcal V\to\mathbb R^d$，并定义
\[
 R_v=f_v-a_v,\qquad I_\nu(f)=I_\nu(a)+I_\nu(R).
\]
其中 $a:\mathcal V\to\mathbb R^d$ 的积分 $I_\nu(a)$ 可以精确计算。标量情形取 $d=1$。令 $A\sim\rho$，假设
\[
 \nu(v)\lVert R_v\rVert>0\ \Longrightarrow\ \rho(v)>0.
\]
这个支持条件只要求覆盖非零残差。定义残差修正估计量
\begin{equation}
 \widehat I_{\nu,\rho,a}(f)
 =I_\nu(a)+\frac{\nu(A)}{\rho(A)}R_A.
 \label{eq:kl-master-estimator}
\end{equation}
比率只在实际抽到的动作上计算。$\rho(v)=0$ 且残差贡献为零的动作无需定义该比率。
\end{definition}

该定义中，$\nu$ 决定要计算哪个积分，$\rho$ 决定从哪里取得随机动作，$a$ 决定哪些贡献已经被确定性地计算。三个选择应当分别记录。特别地，参考分布 $q$ 出现在 $r_v$ 内，并不意味着积分测度或抽样测度就是 $q$。

本节条件于当前参数，把这些函数当作词表上的确定值。估计梯度时直接代入向量函数 $f_v=r_vs_v$ 或 $-s_v$；对一个无偏数值估计量再次自动微分，需要另外验证第~\ref{sec:kl-autodiff} 节的规则。

\begin{theorem}[残差修正的均值、协方差与均方误差]
\label{thm:kl-master-moments}
记 $m_R=I_\nu(R)$。在定义~\ref{def:kl-residual-integral} 的支持条件下，式~\eqref{eq:kl-master-estimator} 无偏，并且
\begin{align}
 \operatorname{Cov}(\widehat I_{\nu,\rho,a}(f))
 &=\sum_{v:\rho(v)>0}\frac{\nu(v)^2}{\rho(v)}R_vR_v^\top-m_Rm_R^\top,
 \label{eq:kl-master-cov}\\
 \Var(\widehat I_{\nu,\rho,a}(f))
 &=\sum_{v:\rho(v)>0}\frac{\nu(v)^2}{\rho(v)}\lVert R_v\rVert^2-\lVert m_R\rVert^2.
 \label{eq:kl-master-trace}
\end{align}
其相对于 $I_\nu(f)$ 的均方误差等于式~\eqref{eq:kl-master-trace}。$n$ 个独立同分布动作对应的样本均值，其协方差和均方误差分别为上述结果的 $1/n$。
\end{theorem}

\begin{proof}
先逐项计算随机修正项的期望。支持条件保证被略去的动作对 $I_\nu(R)$ 没有贡献，因此
\[
 \E_\rho\!\left[\frac{\nu(A)}{\rho(A)}R_A\right]
 =\sum_{v:\rho(v)>0}\nu(v)R_v=I_\nu(R).
\]
加上 $I_\nu(a)$，得到 $I_\nu(a+R)=I_\nu(f)$。确定项不影响协方差；令 $Y=\nu(A)R_A/\rho(A)$，有
\[
 \E[YY^\top]
 =\sum_{v:\rho(v)>0}\frac{\nu(v)^2}{\rho(v)}R_vR_v^\top,
 \qquad \E[Y]=m_R.
\]
应用 $\operatorname{Cov}(Y)=\E[YY^\top]-\E[Y]\E[Y]^\top$ 得到矩阵公式。对矩阵取迹，并使用 $\operatorname{tr}(zz^\top)=\lVert z\rVert^2$，得到迹方差。最后，无偏性给出
\[
 \E\lVert\widehat I-I_\nu(f)\rVert^2
 =\operatorname{tr}\operatorname{Cov}(\widehat I).
\]
独立样本之间的交叉协方差为零，因而平均后的协方差除以 $n$。
\end{proof}

\subsection{精确项误差、错误权重与缺失支持}

残差公式同时给出实现偏差的核对方式。令 $\mathcal D=\{v:\rho(v)>0\}$，允许实际使用的权重 $\widetilde w(v)$ 偏离 $\nu(v)/\rho(v)$，并允许确定积分项存在固定误差 $\delta\in\mathbb R^d$。实际输出为
\[
 \widetilde I=I_\nu(a)+\delta+\widetilde w(A)R_A.
\]
这里先条件于所有已冻结的缓存和控制参数；随机缓存的处理在本节末给出。

\begin{theorem}[三种误差来源的统一偏差]
\label{thm:kl-master-bias}
上述输出相对于 $I_\nu(f)$ 的偏差恰为
\begin{equation}
 b_I=\delta
 +\sum_{v\in\mathcal D}\bigl(\rho(v)\widetilde w(v)-\nu(v)\bigr)R_v
 -\sum_{v\notin\mathcal D}\nu(v)R_v.
 \label{eq:kl-master-bias}
\end{equation}
记 $\widetilde m_R=\sum_{v\in\mathcal D}\rho(v)\widetilde w(v)R_v$，则
\begin{align}
 \operatorname{Cov}(\widetilde I)
 &=\sum_{v\in\mathcal D}\rho(v)\widetilde w(v)^2R_vR_v^\top
 -\widetilde m_R\widetilde m_R^\top,\\
 \E\lVert\widetilde I-I_\nu(f)\rVert^2
 &=\operatorname{tr}\operatorname{Cov}(\widetilde I)+\lVert b_I\rVert^2.
 \label{eq:kl-master-mse}
\end{align}
\end{theorem}

\begin{proof}
由 $I_\nu(f)=I_\nu(a)+\sum_v\nu(v)R_v$，直接相减得
\[
 \E[\widetilde I]-I_\nu(f)
 =\delta+\sum_{v\in\mathcal D}\rho(v)\widetilde w(v)R_v-\sum_v\nu(v)R_v.
\]
把最后一个求和分成 $\mathcal D$ 及其补集，就得到式~\eqref{eq:kl-master-bias}。协方差由随机项的一阶、二阶矩计算。均方误差展开为
\[
 \E\lVert(\widetilde I-\E\widetilde I)+b_I\rVert^2
 =\E\lVert\widetilde I-\E\widetilde I\rVert^2+\lVert b_I\rVert^2,
\]
其中交叉项因 $\E[\widetilde I-\E\widetilde I]=0$ 消失。
\end{proof}

精确积分误差、重要性权重误差和缺失支持具有不同来源，却都在同一个残差上累积。若近似函数已经精确覆盖某个动作，即 $R_v=0$，该动作的权重误差或支持缺失不会产生偏差。反之，即使头部计算完全精确，缺失尾部的残差均值仍会留下系统偏差。无偏性的严格条件是式~\eqref{eq:kl-master-bias} 为零；三个误差项也能相互抵消，因此分别为零是便于验证的充分条件。

\subsection{已知积分控制函数的最优线性组合}

现在固定 $\nu$ 与 $\rho$，假设 $\rho$ 覆盖 $\nu$ 的支持。取 $m$ 个向量控制函数 $h_j:\mathcal V\to\mathbb R^d$，每个积分 $I_\nu(h_j)$ 已知。标量系数 $c_j$ 组成 $c\in\mathbb R^m$，并令
\[
 a_{c,v}=\sum_{j=1}^m c_jh_{j,v},\qquad
 w(v)=\frac{\nu(v)}{\rho(v)},\qquad X=w(A)f_A,
\]
\[
 U_j=w(A)h_{j,A}-I_\nu(h_j),\qquad
 X_0=X-I_\nu(f).
\]
于是 $\E U_j=0$，且主估计量可以逐样本写成
\[
 \widehat I_c=X-\sum_jc_jU_j.
\]
这一步把选择近似函数转化为在零均值随机函数空间中投影。

\begin{theorem}[Gram 矩阵、最优控制与精确超额方差]
\label{thm:kl-master-gram}
定义对称半正定矩阵 $\Gamma$、向量 $\zeta$ 和基准方差 $V_0$：
\[
 \Gamma_{ij}=\E_\rho\langle U_i,U_j\rangle,\qquad
 \zeta_j=\E_\rho\langle U_j,X_0\rangle,\qquad
 V_0=\E_\rho\lVert X_0\rVert^2.
\]
令 $\Gamma^\dagger$ 为 Moore--Penrose 伪逆，则迹方差最优的系数集合为
\[
 c^*=\Gamma^\dagger\zeta+\ker\Gamma.
\]
最优值及任意系数相对于最优值的代价分别为
\begin{align}
 V_{\min}&=V_0-\zeta^\top\Gamma^\dagger\zeta,\\
 \Var(\widehat I_c)-V_{\min}
 &=(c-\Gamma^\dagger\zeta)^\top\Gamma(c-\Gamma^\dagger\zeta).
 \label{eq:kl-master-gram-regret}
\end{align}
这里最小化的是协方差迹；一般向量控制下不同时最小化所有方向的方差。
\end{theorem}

\begin{proof}
估计量的中心化形式为 $X_0-\sum_jc_jU_j$，故展开平方得到
\[
 \Var(\widehat I_c)=V_0-2c^\top\zeta+c^\top\Gamma c.
\]
对任意 $z\in\ker\Gamma$，有
\[
 0=z^\top\Gamma z=\E\left\lVert\sum_jz_jU_j\right\rVert^2.
\]
因此 $\sum_jz_jU_j=0$ 几乎处处成立，进而
\[
 z^\top\zeta=\E\left\langle\sum_jz_jU_j,X_0\right\rangle=0.
\]
这证明 $\zeta\perp\ker\Gamma$，即 $\zeta\in\operatorname{range}\Gamma$。令 $c_0=\Gamma^\dagger\zeta$，便有 $\Gamma c_0=\zeta$。对二次式配方得
\[
 V_0-2c^\top\zeta+c^\top\Gamma c
 =V_0-c_0^\top\Gamma c_0+(c-c_0)^\top\Gamma(c-c_0).
\]
由伪逆恒等式，$c_0^\top\Gamma c_0=\zeta^\top\Gamma^\dagger\zeta$。最后一个二次型非负，且恰在 $c-c_0\in\ker\Gamma$ 时为零，得到全部结论。
\end{proof}

单控制的协方差除以方差公式是该定理在 $m=1$ 时的特例；多个相关控制必须使用整个 Gram 矩阵。控制之间的相关项决定一个控制能否补充另一个控制已经消去的波动。增加控制函数使可选空间包含原来的空间，因此最优迹方差只能下降或保持不变；实际使用不准确的系数，其代价由式~\eqref{eq:kl-master-gram-regret} 精确给出。

\subsection{固定残差时的最优抽样分布}

控制函数改变残差，抽样分布改变残差的二阶矩权重。因此控制优化与抽样优化属于同一方差问题的两个坐标。

\begin{theorem}[残差范数决定迹方差最优抽样]
\label{thm:kl-master-proposal}
固定 $a$，令 $L_R=\sum_v\nu(v)\lVert R_v\rVert$。若 $L_R>0$，则主估计量的最优抽样分布为
\[
 \rho^*(v)=\frac{\nu(v)\lVert R_v\rVert}{L_R},
\]
其最小迹方差为
\[
 \Var_{\min}=L_R^2-\lVert I_\nu(R)\rVert^2.
\]
若 $L_R=0$，则 $I_\nu(a)=I_\nu(f)$，直接返回确定积分项便有零方差。
\end{theorem}

\begin{proof}
令 $\mathcal R=\{v:\nu(v)\lVert R_v\rVert>0\}$。对任何满足残差支持条件的 $\rho$，由 Cauchy--Schwarz 不等式，
\begin{align*}
 L_R^2
 &=\left(\sum_{v\in\mathcal R}\frac{\nu(v)\lVert R_v\rVert}{\sqrt{\rho(v)}}\sqrt{\rho(v)}\right)^2\\
 &\le\left(\sum_{v\in\mathcal R}\frac{\nu(v)^2\lVert R_v\rVert^2}{\rho(v)}\right)
       \left(\sum_{v\in\mathcal R}\rho(v)\right)\\
 &\le\sum_{v\in\mathcal R}\frac{\nu(v)^2\lVert R_v\rVert^2}{\rho(v)}.
\end{align*}
取 $\rho=\rho^*$ 时，全部质量放在 $\mathcal R$，且两向量成比例，两个不等式均取等。代入式~\eqref{eq:kl-master-trace} 得到最优值。若 $L_R=0$，则 $\nu(v)R_v=0$ 对所有动作成立，目标已被精确积分项覆盖。
\end{proof}

在标量情形，最优二阶矩由 $\lvert R\rvert$ 决定；在向量情形，由欧氏范数决定。即使取最优抽样，残差方向随动作变化仍会产生方差。零方差要求三角不等式取等，也就是所有非零残差向量沿同一个有向射线。若额外要求 $\rho$ 对整个词表严格为正，而最优解对零残差动作赋零质量，则上述值是可逼近的下确界。

改变 $a$ 会同时改变 $\Gamma$、$\zeta$ 和 $\rho^*$；改变 $\rho$ 也会改变控制系数的最优解。因而在某个采样分布上拟合出的最优控制系数，不能直接声明在另一个采样分布上仍然最优。

这些最优值是总体矩已知时的 oracle 结论。稀疏 teacher payload 不提供所有未知尾部的 $r_v$，因此一般不能直接构造 $\rho^*$ 或精确计算最优 Gram 系数。它们用于给出可达到的下界、判断控制特征是否有用，以及设计独立样本上的近似拟合；可计算算法必须另外说明这些量如何取得。

\subsection{分区后统一分离事件内部与事件之间的方差}

\begin{theorem}[on-policy 残差的分区协方差]
\label{thm:kl-master-partition}
令 $\rho=\nu$，$\{C_\ell\}_{\ell=1}^L$ 为词表分区，忽略零质量单元。记
\[
 \omega_\ell=\nu(C_\ell),\quad
 \mu_\ell=\E_\nu[R_A\mid A\in C_\ell],\quad
 \Sigma_\ell=\operatorname{Cov}_\nu(R_A\mid A\in C_\ell),\quad
 \bar\mu=\sum_\ell\omega_\ell\mu_\ell.
\]
则
\begin{equation}
 \operatorname{Cov}(\widehat I_{\nu,\nu,a}(f))
 =\sum_\ell\omega_\ell\Sigma_\ell
  +\sum_\ell\omega_\ell(\mu_\ell-\bar\mu)(\mu_\ell-\bar\mu)^\top.
 \label{eq:kl-master-partition}
\end{equation}
\end{theorem}

\begin{proof}
因为重要性比率为一，估计量的协方差就是 $R_A$ 的协方差。条件于 $A\in C_\ell$，写成
\[
 R_A-\bar\mu=(R_A-\mu_\ell)+(\mu_\ell-\bar\mu).
\]
外积展开后，第一项的条件期望为 $\Sigma_\ell$，第二项为确定外积。两个交叉项均含 $\E[R_A-\mu_\ell\mid A\in C_\ell]=0$，故为零。最后按单元概率 $\omega_\ell$ 求和即得结论。
\end{proof}

该公式使头尾控制的作用可以逐项识别。在 $\nu=\rho=p$ 的例子中，取头部 $S$ 与尾部 $S^c$，局部记 $P=p(S)$、$\tau_p=1-P$、$J_A=\mathbf1\{A\notin S\}$。令 $a_v=f_v$ 在头部成立，且 $a_v=b$ 在尾部成立。记尾部条件均值为 $\mu_f$、条件协方差为 $\Sigma_f$，则
\[
 I_p(a)=\sum_{v\in S}p(v)f_v+\tau_pb,\qquad
 R_A=J_A(f_A-b),
\]
\begin{equation}
 \operatorname{Cov}(\widehat I)
 =\tau_p\Sigma_f+P\tau_p(\mu_f-b)(\mu_f-b)^\top.
 \label{eq:kl-master-head-tail}
\end{equation}
减去常向量 $b$ 不改变尾部内部协方差，只改变头尾事件之间的条件均值差。取 $b=\mu_f$ 消去事件间项；要进一步降低事件内部项，需要尾部随动作变化的控制函数。数值目标与梯度目标只需分别代入 $f_v=r_v$ 与 $f_v=r_vs_v$，全部计算便沿用同一分区公式。

\subsection{残差均值、残差波动与随机控制的条件化}

若仅返回确定近似 $I_\nu(a)$，其偏差为 $-I_\nu(R)$，方差为零。若补上正确的残差抽样，偏差变为零，留下的是加权残差的中心二阶矩。由此，近似误差的均值与抽样误差的波动必须分别分析：残差均值为零不保证方差小，残差几乎为常向量也不保证确定近似无偏。

控制函数、系数和采样分布可以由历史数据学习。令 $\mathcal F$ 包含拟合数据和冻结缓存；若条件于 $\mathcal F$，这些对象已固定，且新动作满足 $A\mid\mathcal F\sim\rho_{\mathcal F}$，前述定理均可逐条件应用。若目标 $I_\nu(f)$ 本身固定且每个条件估计量无偏，则
\[
 \E[\widehat I\mid\mathcal F]=I_\nu(f),\qquad
 \operatorname{Cov}(\widehat I)=\E[\operatorname{Cov}(\widehat I\mid\mathcal F)].
\]
若条件偏差为随机向量 $b_{\mathcal F}$，全协方差公式则给出
\[
 \operatorname{Cov}(\widehat I)
 =\E[\operatorname{Cov}(\widehat I\mid\mathcal F)]+\operatorname{Cov}(b_{\mathcal F}),
 \qquad \operatorname{Bias}(\widehat I)=\E[b_{\mathcal F}].
\]
同一动作同时参与控制拟合与残差评估时，条件抽样关系需要重新证明。例如 $\E U=0$ 一般不能推出 $\E[c(A)U]=0$。独立校准集或交叉拟合（cross-fitting）可以使每个评估动作独立于用于构造其控制函数的训练数据，从而恢复逐条件无偏性；交叉拟合各折之间仍能相关，总方差不能直接按独立样本数相除。
